\subsection{Deformación armónica ponderada y selección de la rama superestable}
\label{mab:sec:extension-dodecafasica-ponderada-armonica}

El perfil uniforme posee una extensión paramétrica que conserva las doce
fases y separa la deformación armónica del parámetro dinámico. Para
\(1\le m\le12\), sean

\begin{equation}
 \phi_m=\frac{(30m-10)\pi}{180},
 \qquad
 p(\phi)=12\cos(3\phi)-\cos(4\phi).
 \label{mab:eq:fases-y-deformacion-armonica-armonizada}
\end{equation}

Para todo \(\eta\) real tal que los pesos permanezcan positivos, definimos

\begin{equation}
 w_m(\eta)=\frac{1+\eta p(\phi_m)}{12}.
 \label{mab:eq:pesos-dodecafasicos-armonizados}
\end{equation}

Las sumas de los armónicos tercero y cuarto sobre las doce fases se anulan;
por tanto,

\[
 \sum_{m=1}^{12}w_m(\eta)=1.
\]

Sea

\begin{equation}
 s_\eta^2=\frac{2}{\sum_{m=1}^{12}w_m(\eta)\phi_m^2},
 \qquad
 u_\eta(x)=1-\sum_{m=1}^{12}w_m(\eta)
 \cos(s_\eta\phi_m x).
 \label{mab:eq:familia-ponderada-armonizada}
\end{equation}

La expansión del coseno y la definición de \(s_\eta\) dan

\[
 u_\eta(x)=x^2+O(x^4)\qquad(x\to0).
\]

Por consiguiente, \(u_\eta\) es real--analítica, par y conserva un punto
crítico cuadrático en el origen. Introducimos el parámetro dinámico mediante

\begin{equation}
 f_{\lambda,\eta}(x)=1-\lambda u_\eta(x).
 \label{mab:eq:familia-feigenbaum-ponderada-armonizada}
\end{equation}

Un parámetro superestable de periodo exacto \(2^n\) satisface

\[
 f_{\lambda,\eta}^{\,2^n}(0)=0,
 \qquad
 f_{\lambda,\eta}^{\,2^k}(0)\ne0\quad(0\le k<n).
\]

La ecuación puede poseer varias ramas. Se fija una rama
\(\mathscr B\) recorriendo, inmediatamente después del parámetro del
nivel anterior, una malla adaptativa; se toma el primer intervalo con cambio
de signo y se aplica bisección. La notación
\(\lambda_n^{\mathscr B}(\eta)\) incorpora este selector y no atribuye
unicidad global a la ecuación superestable. El primer nivel se determina
exactamente:

\[
 \lambda_1^{\mathscr B}(\eta)=\frac1{u_\eta(1)},
\]

si \(u_\eta(1)>0\), porque
\(f_{\lambda,\eta}^{\,2}(0)=f_{\lambda,\eta}(1)=0\).
Dados tres niveles consecutivos y

\[
 x_n^{\mathscr B}(\eta)
 :=f_{\lambda_n^{\mathscr B}(\eta),\eta}^{\,2^{n-1}}(0),
\]

se forman

\begin{equation}
 \delta_n^{\mathscr B}(\eta)
 =\frac{\lambda_{n-1}^{\mathscr B}(\eta)-
         \lambda_{n-2}^{\mathscr B}(\eta)}
        {\lambda_n^{\mathscr B}(\eta)-
         \lambda_{n-1}^{\mathscr B}(\eta)},
 \qquad
 \alpha_n^{\mathscr B}(\eta)
 =\frac{|x_{n-1}^{\mathscr B}(\eta)|}
        {|x_n^{\mathscr B}(\eta)|}.
 \label{mab:eq:aproximantes-ponderados-armonizados}
\end{equation}

Para \(\eta=0\) se recupera exactamente el perfil uniforme. La
extensión ponderada introduce un parámetro transversal: permite estudiar
la transversalidad de la familia respecto de la variedad estable del
renormalizador sin convertir una coincidencia finita en una prueba de
universalidad.

\begin{remark}[Control negativo]
La extensión queda refutada si los pesos pierden positividad en el dominio
declarado, si la normalización deja de producir coeficiente cuadrático
unitario, si el selector no excluye periodos divisores o si una rama se elige
usando como entrada los valores universales terminales.
\end{remark}
